\documentclass[exercice]{thedocument}%
\startdatakeys
\source{}
\cycle{}
\competencesDuSocle{}
\connaissancesRequises{}
\competencesTravaillees{}
\enddatakeys
\begin{document}%
\begin{minipage}[t]{.5\linewidth}
    Dans ce problème, on lance deux dés de couleurs différentes. Les dés sont
    équilibrés et les faces sont numérotées de 1 à 6. On s'intéresse à la somme
    des valeurs obtenues par les dés.\par\noindent
    {\bfseries Partie 1 : }On lance 25 fois les deux dés et on note les valeurs dans un
    tableur.\par\noindent Les résultats sont représentés dans le tableau ci-contre.
    \par\noindent La colonne A indique le numéro de l'expérience. Les colonnes B
    et C donnent les valeurs des dés. La sommme des deux dés est calculée dans la
    colonne D.
    \begin{enumerate}
        \item La somme peut-elle être égale à 1~? Justifier.
        \item La somme 12 n'apparaît pas dans ce tableau. Est-il toutefois possible
        de l'obtenir~? Justifier.
        \item Pour le 11$^e$ lancer des deu dés, quelle formule a-t-on marquée
        dans la cellule D12 pour obtenir le résultat donné par l'ordinateur~?
        \item Dans cette expérience, combien de fois obtient-on la somme 7~?
        \par\noindent En déduire la fréquence de cette somme en pourcentage.
        \item Quelle est la médiane de cette série de sommes (colonne D)~?
        \item Tracer le diagramme en bâtons de la sériedes sommes obtenues
        (colonne D).
    \end{enumerate}
\end{minipage}\hspace{0.1\linewidth}
\begin{minipage}[t]{.39\linewidth}
    \centering\pdf[scale=1]{Tableau_stats_GEN2010POL3}\par
\end{minipage}\bigskip\par\noindent
{\bfseries Partie 2 : }On fait une simulation de $1\,000$ expériences avec un tableur.
Les résultats sont représentés dans le diagramme en bâtons suivant.
\centering\pdf[scale=0.8]{diagramme_batons_GEN2010POL3}\par
\begin{enumerate}
    \item Quelles sont les deux sommes les moins fréquentes~?
    \item Paul, un élève de troisième joue avec Jacques son petit frère de CM2.
    Chacun choisit une somme à obteniravec 2 dés. Paul prend la somme 9
    et Jacques la somme 3.\vskip10pt Expliquez pourquoi Paul a plus de chances
    de gagner que son petit frère.
    \item Quel est, pour cette simulation, le nombre de lancers qui donne la
    somme 7~? En déduire la fréquence en pourcentage représentée par ces lancers.
    \item {\bfseries Compléter le tableau suivant sur cette feuille} et trouver les
    différentes possibilités d'obtenir une somme égale à 7 avec deux dés. Calculer
    la probabilité d'obtenir cette somme.
    \par\noindent\centering\pdf[scale=1]{tableau_GEN2010POL3}\par\noindent
    \item Que peut-on dire de la valeur de la fréquence obtenue à la question 3
    et de celle de la probabilité obtenue à la question 4~? Proposer une
    explication.
\end{enumerate}
\end{document}
